%!TEX TS-program = pdflatex
\documentclass[aspectratio=169,10pt]{beamer}
\usepackage[T2A]{fontenc}
\usepackage[utf8]{inputenc}
\usepackage[english,russian]{babel}
\usepackage{cmap}
\usepackage{amsmath,amssymb,mathtools,booktabs,array,etoolbox,PTSans,graphicx}
\graphicspath{{figures/}}
\renewcommand{\familydefault}{\sfdefault}
\renewcommand{\rmdefault}{\sfdefault}
\newbool{russian}\booltrue{russian}
\usepackage{HSE-theme/beamerthemeHSE}
\hypersetup{bookmarks=false}
\setbeamertemplate{blocks}[rounded][shadow=false]
\setbeamerfont{frametitle}{size=\Large}
\setbeamerfont{framesubtitle}{size=\small}
\setbeamerfont{block title}{size=\normalsize}
\setbeamerfont{block body}{size=\small}
\setlength{\abovedisplayskip}{5pt}
\setlength{\belowdisplayskip}{5pt}
\newcommand{\R}{\mathbb{R}}
\newcommand{\norm}[1]{\left\lVert #1\right\rVert}
\newcommand{\grad}{\nabla}
\newcommand{\figwide}[1]{\begin{center}\includegraphics[height=.50\textheight,keepaspectratio]{#1}\end{center}}
\newcommand{\figshort}[1]{\begin{center}\includegraphics[height=.41\textheight,keepaspectratio]{#1}\end{center}}
\title[Методы оптимизации]{Методы оптимизации\\в машинном обучении}
\subtitle{Лекция 10. Предобусловливание и геометрия методов оптимизации}
\author[Н. Лукьяненко]{Никита Лукьяненко}
\institute{Совместная программа СМОЛ ГУ и ФКН НИУ ВШЭ}
\date{2026/27 учебный год}
\begin{document}

% Блок 1. Мотивация, 7 минут. Слайды 1--6.
% 1
\frame[plain]{\titlepage}

% 2
\begin{frame}{От градиентного спуска к предобусловливанию}
\small
В лекциях 4 и 6 мы связали скорость градиентного спуска с различием кривизны по направлениям. В лекции 9 метод Ньютона исправлял это различие при помощи гессиана.
\[
x_{k+1}^{GD}=x_k-\alpha_k\grad f(x_k),\qquad
x_{k+1}^{N}=x_k-\alpha_k H(x_k)^{-1}\grad f(x_k).
\]
\begin{block}{Вопрос этой лекции}
Можно ли использовать \textbf{простую} матрицу $P\succ0$, которая отражает масштаб направлений и обходится дешевле точного гессиана?
\end{block}
\[
\boxed{x_{k+1}=x_k-\alpha_kP^{-1}\grad f(x_k).}
\]
\end{frame}

% 3
\begin{frame}{Что мы уже знаем об обусловленности}
\small
Для квадратичной функции $f(x)=\tfrac12x^\top Ax-b^\top x$ с $A\succ0$ линии уровня вытянуты, если собственные значения $A$ сильно различаются.
\[
\kappa(A)=\frac{\lambda_{\max}(A)}{\lambda_{\min}(A)}.
\]
При шаге $1/\lambda_{\max}(A)$ малая кривизна замедляет градиентный спуск.
\begin{alertblock}{Сегодняшний акцент}
Число обусловленности уже изучалось. Теперь будем менять \textbf{геометрию алгоритма}, чтобы изменить спектр, который определяет его скорость.
\end{alertblock}
\end{frame}

% 4
\begin{frame}{Одна длина шага не подходит всем направлениям}
\small
Рассмотрим $f(x_1,x_2)=\tfrac12(25x_1^2+x_2^2)$. Градиент равен $(25x_1,x_2)^\top$.
Если $\alpha$ подходит направлению $x_1$, то вдоль $x_2$ движение медленное. Увеличение $\alpha$ ради $x_2$ делает движение по $x_1$ неустойчивым.
\begin{block}{Идея}
Разрешить разный масштаб перемещения в разных направлениях:
\[
p=-P^{-1}\grad f(x),\qquad P=\operatorname{diag}(25,1).
\]
На этой задаче $p=-x$ и полный шаг попадает в минимум.
\end{block}
\end{frame}

% 5
\begin{frame}{Цели лекции}
\small
После лекции мы сможем:
\begin{enumerate}
\item вывести предобусловленный градиентный спуск из замены координат;
\item доказать, что при $P\succ0$ его направление убывает;
\item определить допустимый и оптимальный постоянный шаг на квадратичной задаче;
\item связать скорость с собственными значениями $P^{-1/2}AP^{-1/2}$;
\item выбрать простое $P$ и объяснить, когда оно не помогает.
\end{enumerate}
\end{frame}

% 6
\begin{frame}{Маршрут на 80 минут}
\small
\begin{tabular}{@{}p{.71\linewidth}r@{}}
\toprule
Тема & Время\\
\midrule
Мотивация и связь с предыдущими лекциями & 7 мин\\
Замена координат и геометрия & 14 мин\\
Предобусловленный градиентный спуск & 15 мин\\
Спектр и оценки скорости & 18 мин\\
Выбор матрицы $P$ & 11 мин\\
Примеры из машинного обучения & 9 мин\\
Закрепление и переход к следующей лекции & 6 мин\\
\bottomrule
\end{tabular}
\end{frame}

% Блок 2. Замена координат, 14 минут. Слайды 7--16.
% 7
\begin{frame}{Как измерять длину шага}
\footnotesize
Градиентный спуск меняет точку на $h=x_{k+1}-x_k=-\alpha\nabla f(x_k)$.
Обычная длина этого перемещения равна
$\norm h_2=\sqrt{h_1^2+\cdots+h_d^2}$: все координаты считаются в одинаковых единицах.
\par\vspace{.3em}
Но изменение разных координат на единицу может по-разному менять функцию.
Для диагональной $P=\operatorname{diag}(p_1,\ldots,p_d)$ зададим
\[
\norm h_P=\sqrt{h^\top Ph}=\sqrt{p_1h_1^2+\cdots+p_dh_d^2}.
\]
Большое $p_i$ означает: перемещение по координате $i$ считаем более значительным.
\figshort{metric_balls.pdf}
\end{frame}

% 8
\begin{frame}{Простое изменение масштаба}
\small
Для $f(x_1,x_2)=\tfrac12(25x_1^2+x_2^2)$ введём новые координаты
\[
y_1=5x_1,\qquad y_2=x_2.
\]
Тогда $\widetilde f(y)=f(y_1/5,y_2)=\tfrac12(y_1^2+y_2^2)$.
\begin{block}{Геометрический результат}
В координатах $y$ линии уровня становятся окружностями. Обычный градиентный спуск по $y$ одинаково масштабирует оба направления.
\end{block}
\end{frame}

% 9
\begin{frame}{Та же задача в других координатах}
\small
На предыдущем примере мы ввели $y_1=5x_1$, $y_2=x_2$.
Точка $x=(1,2)$ стала записываться как $y=(5,2)$: точка и значение функции не изменились, изменились только числа, которыми её обозначаем.
\par\vspace{.5em}
В общем случае выбираем обратимую матрицу $S$ и записываем $x=Sy$.
Чтобы вычислить значение функции по новым координатам, подставляем $Sy$ вместо $x$:
\[
\widetilde f(y)=f(Sy).
\]
Если $y^*$ минимизирует $\widetilde f$, то $x^*=Sy^*$ минимизирует исходную $f$.
\begin{block}{Зачем это делать}
Алгоритм будет делать шаги по координатам $y$. После перевода этих шагов в $x$ окажется, что шаги по разным направлениям получили разные масштабы.
\end{block}
\end{frame}

% Новая памятка перед правилом цепочки.
% 10
\begin{frame}{Напоминание об обозначениях}
\small
\begin{itemize}
\item $\nabla f(x)$ --- столбец частных производных
$(\partial f/\partial x_1,\ldots,\partial f/\partial x_d)^\top$.
\item $D_vf(x)$ --- производная по направлению $v$:
$D_vf(x)=\left.\frac{d}{dt}f(x+tv)\right|_{t=0}=\nabla f(x)^\top v$.
\item $S^\top$ --- транспонированная матрица. Если в новых координатах мы сдвинулись на $v$, то при $x=Sy$ исходная точка сдвинулась на $Sv$.
\end{itemize}
\begin{block}{Идея следующего вычисления}
Одну и ту же производную по направлению можно найти через $\widetilde f(y)$ или через $f(x)$; сравнение ответов даст формулу для нового градиента.
\end{block}
\end{frame}

% 11
\begin{frame}{Как преобразуется градиент}
\small
Для любого направления $v$ в координатах $y$ соответствующее направление в $x$ равно $Sv$. Поэтому по правилу цепочки
\[
D_v\widetilde f(y)=D_{Sv}f(Sy)
=\grad f(Sy)^\top Sv
=\bigl(S^\top\grad f(Sy)\bigr)^\top v.
\]
Равенство верно при любом $v$, значит стоящие перед $v$ столбцы равны:
\[
\boxed{\grad_y\widetilde f(y)=S^\top\grad_x f(x),\qquad x=Sy.}
\]
Проверка размеров: $S^\top$ и $\nabla_x f$ дают столбец из $d$ компонент, как и $\nabla_y\widetilde f$.
\end{frame}

% 12
\begin{frame}{Как преобразуется гессиан}
\small
Если $S$ постоянно, дифференцируем формулу для градиента ещё раз:
\[
\boxed{\grad_y^2\widetilde f(y)=S^\top H(x)S,\qquad x=Sy.}
\]
Для квадратичной $f(x)=\tfrac12x^\top Ax-b^\top x$ новая матрица кривизны равна $S^\top AS$.
\begin{block}{Следствие}
Правильно выбранный $S$ может приблизить $S^\top AS$ к единичной матрице. Тогда обычный GD в координатах $y$ видит более равномерную кривизну.
\end{block}
\end{frame}

% 13
\begin{frame}{Линии уровня до и после замены координат}
\small
Для $A=\operatorname{diag}(25,1)$ возьмём $S=\operatorname{diag}(1/5,1)$. Тогда $S^\top AS=I$.
\figwide{ellipse_scaling.pdf}
\end{frame}

% 14
\begin{frame}{Поворот и масштабирование действуют по-разному}
\small
Если $S$ ортогональна, то $S^\top AS$ имеет те же собственные значения, что и $A$: мы лишь повернули оси.
\[
S^\top S=I\quad\Longrightarrow\quad\kappa(S^\top AS)=\kappa(A).
\]
Масштабирование способно изменить отношение собственных значений.
\begin{alertblock}{Вывод}
Для ускорения GD недостаточно переименовать координаты. Нужно изменить масштаб вдоль направлений, в которых кривизна различается.
\end{alertblock}
\end{frame}

% 15
\begin{frame}{От замены координат к матрице $P$}
\small
Обычный шаг GD для $\widetilde f$ имеет вид $y_+=y-\alpha\grad_y\widetilde f(y)$.
Используя $x=Sy$ и $\grad_y\widetilde f=S^\top\grad f(x)$, получаем
\[
x_+=Sy_+=x-\alpha SS^\top\grad f(x).
\]
Обозначим $P=(SS^\top)^{-1}\succ0$. Тогда
\[
\boxed{x_+=x-\alpha P^{-1}\grad f(x).}
\]
Это и есть предобусловленный градиентный спуск.
\end{frame}

% 16
\begin{frame}{Выбор $S=P^{-1/2}$}
\small
Для любой симметричной $P\succ0$ существует симметричный корень $P^{1/2}$.
Удобно положить
\[
x=P^{-1/2}y,\qquad y=P^{1/2}x.
\]
Тогда $S=P^{-1/2}$, $SS^\top=P^{-1}$, а гессиан квадратичной задачи в новых координатах равен
\[
\boxed{B=P^{-1/2}AP^{-1/2}.}
\]
Именно спектр $B$ будет определять шаг и скорость. Это главная матрица лекции.
\end{frame}

% Блок 3. Предобусловленный GD, 15 минут. Слайды 17--26.
% 17
\begin{frame}{Предобусловленный градиентный спуск}
\small
Фиксируем симметричную положительно определённую матрицу $P$. На шаге $k$ решаем систему
\[
Pp_k=-\grad f(x_k),\qquad x_{k+1}=x_k+\alpha_kp_k.
\]
Запись $p_k=-P^{-1}g_k$ описывает ответ системы, но саму матрицу $P^{-1}$ хранить не нужно.
Если $P=\operatorname{diag}(p_1,\ldots,p_d)$, то $(p_k)_i=-(g_k)_i/p_i$.
Для более общей $P$ применяют обычный метод решения системы $Pp_k=-g_k$.
\begin{block}{Частные случаи}
$P=I$ даёт обычный градиентный спуск. Если для квадратичной функции взять $P=A$, направление совпадает с ньютоновским.
\end{block}
\end{frame}

% 18
\begin{frame}{Теорема: направление спуска}
\small
\begin{block}{Утверждение}
Пусть $f$ дифференцируема в $x$, $P\succ0$ и $g=\grad f(x)\ne0$.
Тогда $p=-P^{-1}g$ является \textbf{направлением спуска}: существует $\varepsilon>0$, для которого $f(x+\alpha p)<f(x)$ при всех $0<\alpha<\varepsilon$.
\end{block}
\begin{block}{Доказательство}
Матрица $P^{-1}$ положительно определена, поэтому
\[
D_pf(x)=g^\top p=-g^\top P^{-1}g<0.
\]
По определению производной
$f(x+\alpha p)=f(x)+\alpha D_pf(x)+o(\alpha)$.
Для достаточно малого положительного $\alpha$ отрицательное линейное слагаемое определяет знак разности.
\end{block}
\end{frame}

% 19
\begin{frame}{Предобусловливание не меняет стационарные точки}
\small
Так как $P^{-1}$ обратима,
\[
p(x)=0
\quad\Longleftrightarrow\quad
P^{-1}\grad f(x)=0
\quad\Longleftrightarrow\quad
\grad f(x)=0.
\]
Следовательно, предобусловливание меняет шаги алгоритма, но сохраняет все точки, в которых градиент равен нулю.
\begin{alertblock}{Оговорка}
Равенство $\nabla f(x)=0$ ещё не означает минимум. Например,
$f(x_1,x_2)=x_1^2-x_2^2$ имеет нулевой градиент в $(0,0)$, но при движении по оси $x_2$ значение уменьшается. Это седловая точка; любая $P\succ0$ оставит её стационарной.
\end{alertblock}
\end{frame}

% 20
\begin{frame}{Эквивалентность обычному GD в новых координатах}
\small
Пусть каждой точке $x_k$ соответствует $y_k=P^{1/2}x_k$, а
$\widetilde f(y)=f(P^{-1/2}y)$. По правилу цепочки
\[
\grad\widetilde f(y)=P^{-1/2}\grad f(x).
\]
Сделаем обычный шаг GD в $y$ и запишем полученную точку обратно в $x$:
\[
x_+=P^{-1/2}y_+=x-\alpha P^{-1}\grad f(x).
\]
\begin{block}{Смысл}
Такое соответствие верно \textbf{на каждом шаге}. На практике можно обновлять сразу $x_k$ по последней формуле: явно переводить точки в $y$ и обратно не требуется.
\end{block}
\end{frame}

% 21
\begin{frame}{Квадратичная задача: рекурсия ошибки}
\footnotesize
Пусть $f(x)=\tfrac12x^\top Ax-b^\top x$, $A\succ0$, $x^*=A^{-1}b$.
Из $Ax^*=b$ следует $\nabla f(x_k)=A(x_k-x^*)$. Обозначим ошибку $e_k=x_k-x^*$.
Вычитая $x^*$ из формулы шага, получаем
\[
e_{k+1}=(I-\alpha P^{-1}A)e_k.
\]
Матрица $P^{-1}A$ обычно не симметрична. Чтобы анализировать её собственные направления, заменим ошибку на $u_k=P^{1/2}e_k$.
Умножаем предыдущую формулу слева на $P^{1/2}$ и подставляем $e_k=P^{-1/2}u_k$:
\[
\boxed{u_{k+1}=(I-\alpha B)u_k,\quad B=P^{-1/2}AP^{-1/2}\succ0.}
\]
Матрица $B$ симметрична, поскольку $A$ и $P^{-1/2}$ симметричны. Поэтому её можно диагонализовать в ортонормированном базисе.
\end{frame}

% 22
\begin{frame}{Пять шагов из одной начальной точки}
\small
Сравним обычный GD и предобусловленный GD для
$A=\begin{pmatrix}10&2\\2&1\end{pmatrix}$, $P=\operatorname{diag}(10,1)$.
Оба метода начинают с $x_0=(1,3)^\top$.
Для обычного GD берём $\alpha=1/\lambda_{\max}(A)\approx0.096$, для метода с $P$ --- $\alpha=0.3$.
\figwide{path_comparison.pdf}
Слева каждый отрезок --- один шаг; справа показаны значения функции.
Первый шаг обычного GD здесь удачнее, но к пятому шагу метод с $P$ достигает меньшего значения функции.
\end{frame}

% 23
\begin{frame}{Выбор длины шага остаётся необходимым}
\small
После выбора $p_k=-P^{-1}g_k$ ещё нужно выбрать $\alpha_k$.
Можно использовать уже знакомый поиск Армихо:
\[
f(x_k+\alpha_kp_k)
\le f(x_k)+c\alpha_k g_k^\top p_k,\qquad 0<c<1.
\]
Так как $g_k^\top p_k<0$ вне стационарной точки, достаточно малый шаг проходит проверку.
\begin{alertblock}{На практике}
Качество $P$ и выбор $\alpha_k$ решают разные вопросы. Матрица исправляет масштабы направлений, поиск шага контролирует размер перемещения.
\end{alertblock}
\end{frame}

% 24
\begin{frame}{Как оценить уменьшение функции}
\footnotesize
Условие $\nabla^2f(z)\preceq L_PP$ означает:
в любом направлении $v$ кривизна не больше $L_Pv^\top Pv$.
Пусть оно выполняется во всех точках отрезка от $x$ до $x+v$.
Тогда аналог леммы об убывании из лекции 4 даёт
\[
f(x+v)\le f(x)+g^\top v+\tfrac{L_P}{2}v^\top Pv.
\]
Здесь $g=\nabla f(x)$, а $v^\top Pv$ --- квадрат длины $v$ с учётом выбранных весов $P$.
Подставим шаг $v=-\alpha P^{-1}g$:
\[
f(x+v)\le f(x)-\alpha\left(1-\tfrac{\alpha L_P}{2}\right)g^\top P^{-1}g.
\]
При $\alpha=1/L_P$ правая часть заведомо меньше $f(x)$, если $g\ne0$:
\[
\boxed{f(x_+)\le f(x)-\frac{1}{2L_P}g^\top P^{-1}g.}
\]
\end{frame}

% 25
\begin{frame}{Проверка на диагональной квадратичной функции}
\small
Для $f(x_1,x_2)=\tfrac12(25x_1^2+x_2^2)$ возьмём $P=\operatorname{diag}(25,1)$.
Тогда $\grad f(x)=Px$ и
\[
p=-P^{-1}\grad f(x)=-x,\qquad x_+=x+\alpha p=(1-\alpha)x.
\]
При $\alpha=1$ получаем $x_+=0$ из любой начальной точки.
\begin{block}{Вопрос аудитории}
Что будет при $\alpha=2$? Подставьте его в формулу для $x_+$ и проследите первые два шага.
\end{block}
\end{frame}

% 26
\begin{frame}{Проверка на диагональной квадратичной функции}
\small
Для $f(x_1,x_2)=\tfrac12(25x_1^2+x_2^2)$ возьмём $P=\operatorname{diag}(25,1)$.
Тогда $\grad f(x)=Px$ и
\[
p=-P^{-1}\grad f(x)=-x,\qquad x_+=x+\alpha p=(1-\alpha)x.
\]
При $\alpha=1$ получаем $x_+=0$ из любой начальной точки.
\begin{block}{Ответ}
При $\alpha=2$ имеем $x_{k+1}=-x_k$: последовательность точек равна
$x_0,-x_0,x_0,\ldots$.
Расстояние до минимума не уменьшается.
\end{block}
\end{frame}

% Блок 4. Спектр и скорость, 18 минут. Слайды 27--37.
% 27
\begin{frame}{Какой спектр определяет скорость}
\small
В этом разделе предполагаем $A=A^\top\succ0$ и $P=P^\top\succ0$.
Только при этих условиях $B=P^{-1/2}AP^{-1/2}$ гарантированно симметрична и положительно определена:
\[
B^\top=B,\qquad
z^\top Bz=(P^{-1/2}z)^\top A(P^{-1/2}z)>0\quad(z\ne0).
\]
Пусть её собственные значения лежат между
\[
\mu_P=\lambda_{\min}(B)>0,\qquad L_P=\lambda_{\max}(B).
\]
В собственном направлении с $\lambda_i(B)$ ошибка умножается на $1-\alpha\lambda_i(B)$.
\[
\boxed{\kappa_P=\frac{L_P}{\mu_P}}
\]
называем числом обусловленности \textbf{предобусловленной} задачи.
Если $A$ не положительно определена, положительность $B$ из выбора $P$ не следует.
\end{frame}

% 28
\begin{frame}{Обобщённая задача на собственные значения}
\small
Матрица $P^{-1}A$ появилась в рекурсии ошибки:
$e_{k+1}=(I-\alpha P^{-1}A)e_k$ (слайд о квадратичной задаче).
Ведь $\nabla f(x_k)=Ae_k$, а направление равно $-P^{-1}\nabla f(x_k)$.
\par\vspace{.3em}
$P^{-1}A$ не обязана быть симметричной. Однако она подобна симметричной матрице $B$:
\[
P^{1/2}(P^{-1}A)P^{-1/2}=P^{-1/2}AP^{-1/2}=B.
\]
Поэтому обе матрицы имеют одинаковые положительные собственные значения.
\begin{block}{Другая запись}
Уравнение $P^{-1}Av=\lambda v$ можно записать как $Av=\lambda Pv$.
Это обобщённая задача на собственные значения; явное вычисление $P^{-1}$ для её записи не нужно.
\end{block}
\end{frame}

% 29
\begin{frame}{Теорема: допустимый постоянный шаг}
\small
\begin{block}{Утверждение для квадратичной функции}
Для любого начального $x_0$ итерации предобусловленного GD сходятся к $x^*$ тогда и только тогда, когда
\[
\boxed{0<\alpha<\frac{2}{L_P}.}
\]
\end{block}
\begin{block}{Доказательство}
Повторяем доказательство сходимости GD на квадратичной функции из лекции 4:
там ошибку раскладывали по собственным векторам $A$, здесь --- по собственным векторам $B$.
Если $u_0=\sum_i c_iq_i$, то
\[
u_k=\sum_i c_i(1-\alpha\lambda_i)^kq_i.
\]
Для сходимости из \emph{любой} начальной точки нужно
$|1-\alpha\lambda_i|<1$ при каждом $i$, то есть $0<\alpha\lambda_i<2$.
Наибольшее $\lambda_i=L_P$ даёт указанную границу.
\end{block}
\end{frame}

% 30
\begin{frame}{Оптимальный постоянный шаг на квадратичной задаче}
\footnotesize
Худший множитель по собственным направлениям равен
\[
\rho(\alpha)=\max_{\lambda\in[\mu_P,L_P]}|1-\alpha\lambda|
=\max\{|1-\alpha\mu_P|,\ |1-\alpha L_P|\}.
\]
Минимум достигается, когда значения на концах отрезка равны по модулю:
\[
1-\alpha\mu_P=\alpha L_P-1.
\]
Отсюда
\[
\boxed{\alpha_*=\frac{2}{\mu_P+L_P},\qquad
\rho_*=\frac{L_P-\mu_P}{L_P+\mu_P}
=\frac{\kappa_P-1}{\kappa_P+1}.}
\]
\end{frame}

% 31
\begin{frame}{Оценка сходимости квадратичной функции}
\small
\begin{block}{Теорема}
При $\alpha=\alpha_*$ для всех $k\ge0$:
\[
\norm{x_k-x^*}_P\le\rho_*^k\norm{x_0-x^*}_P,
\qquad
f(x_k)-f^*\le\rho_*^{2k}\bigl(f(x_0)-f^*\bigr).
\]
\end{block}
Обозначения: $\norm e_P=\sqrt{e^\top Pe}$ и
$\rho_*=(\kappa_P-1)/(\kappa_P+1)$.
Докажем обе оценки на следующем слайде.
\end{frame}

% 32
\begin{frame}{Доказательство оценки на квадратичной функции}
\footnotesize
Разложим $u_0=P^{1/2}(x_0-x^*)=\sum_i c_iq_i$ по ортонормированным собственным векторам $B$.
Тогда $u_k=\sum_i c_i(1-\alpha_*\lambda_i)^kq_i$ и
$|1-\alpha_*\lambda_i|\le\rho_*$.
\par\vspace{.4em}
Поскольку $\norm{x_k-x^*}_P=\norm{u_k}_2$,
\[
\norm{u_k}_2^2=\sum_i c_i^2(1-\alpha_*\lambda_i)^{2k}
\le\rho_*^{2k}\sum_i c_i^2=\rho_*^{2k}\norm{u_0}_2^2.
\]
Для квадратичной функции $f(x_k)-f^*=\tfrac12u_k^\top Bu$, поэтому
\[
f(x_k)-f^*
=\tfrac12\sum_i\lambda_i c_i^2(1-\alpha_*\lambda_i)^{2k}
\le\rho_*^{2k}\,\tfrac12\sum_i\lambda_i c_i^2
=\rho_*^{2k}(f(x_0)-f^*).
\]
\end{frame}
% 33
\begin{frame}{Как зависит худший множитель от шага}
\small
На примере с $\mu_P=0.1$ и $L_P=1.9$ оптимальный шаг равен $1$.
\figwide{step_spectrum.pdf}
Слишком большой шаг увеличивает ошибку в направлении максимальной кривизны.
\end{frame}

% 34
\begin{frame}{Теорема для неквадратичной функции}
\footnotesize
\begin{block}{Глобальная оценка при фиксированном $P$}
Пусть $f\in C^2(\R^d)$, $P=P^\top\succ0$ и для всех $x$ выполнено
\[
\mu_PP\preceq\grad^2f(x)\preceq L_PP,\qquad
0<\mu_P\le L_P.
\]
Тогда $f$ имеет единственный минимум $x^*$. Для
$x_{k+1}=x_k-\frac1{L_P}P^{-1}\grad f(x_k)$:
\[
\boxed{f(x_k)-f^*\le
\left(1-\frac{\mu_P}{L_P}\right)^k
\bigl(f(x_0)-f^*\bigr).}
\]
\end{block}
Условие сравнивает кривизну $f$ с одной и той же матрицей $P$ во всей $\R^d$.
Это аналог оценки из лекции 6 после замены координат $y=P^{1/2}x$.
\end{frame}

% 35
\begin{frame}{Доказательство: что даёт верхняя граница}
\footnotesize
Положим $g=\nabla f(x)$ и $s=-L_P^{-1}P^{-1}g$.
На отрезке от $x$ до $x+s$ верхняя граница гессиана даёт
\[
f(x+s)\le f(x)+g^\top s+\frac{L_P}{2}s^\top Ps.
\]
Подставляя $s$, получаем
\[
f(x_+)\le f(x)-\frac{1}{2L_P}g^\top P^{-1}g.
\]
Это количественная оценка убывания за один шаг.
Чтобы связать её с расстоянием до оптимального значения $f^*$, используем нижнюю границу гессиана на следующем слайде.
\end{frame}

% 36
\begin{frame}{Доказательство: что даёт нижняя граница}
\footnotesize
Из $\nabla^2f(z)\succeq\mu_PP$ следует для любой точки $y$
\[
f(y)\ge f(x)+g^\top(y-x)+\tfrac{\mu_P}{2}\norm{y-x}_P^2.
\]
Минимум правой части по $h=y-x$ достигается при
$h=-\mu_P^{-1}P^{-1}g$ и равен
$f(x)-\frac{1}{2\mu_P}g^\top P^{-1}g$.
Поскольку $f^*=\min_y f(y)$, имеем
\[
f(x)-f^*\le\frac{1}{2\mu_P}g^\top P^{-1}g.
\]
Совмещаем это с оценкой предыдущего слайда:
\[
f(x_+)-f^*\le(1-\mu_P/L_P)\bigl(f(x)-f^*\bigr).
\]
Повторение неравенства $k$ раз доказывает теорему.
\end{frame}

% 37
\begin{frame}{Один старт, две скорости}
\small
Для $A=\begin{pmatrix}100&9\\9&1\end{pmatrix}$ исходное $\kappa(A)\approx535$, а после диагонального масштабирования $\kappa_P=19$.
\figwide{convergence_comparison.pdf}
В этом примере матрица $P$ существенно уменьшает разброс кривизны и число шагов до заданной точности.
\end{frame}

% Блок 5. Выбор матрицы P, 11 минут. Слайды 38--45.
% 38
\begin{frame}{Идеальный случай $P=A$}
\small
Для квадратичной функции при $P=A$ получаем $B=A^{-1/2}AA^{-1/2}=I$ и $\kappa_P=1$.
\[
x_+=x-\alpha A^{-1}(Ax-b).
\]
При $\alpha=1$ это $x_+=A^{-1}b=x^*$ за один шаг --- результат лекции 9.
\begin{alertblock}{Цена точности}
Если $A$ велика и плотна, решение системы с $A$ может быть слишком дорогим. Поэтому мы ищем $P$, которую проще применять и которая \emph{достаточно} выравнивает спектр.
\end{alertblock}
\end{frame}

% 39
\begin{frame}{Диагональное предобусловливание}
\small
Простой выбор для квадратичной задачи:
\[
\boxed{P=\operatorname{diag}(A_{11},\ldots,A_{dd}).}
\]
Такой $P$ положительно определён, поскольку у $A\succ0$ все диагональные элементы положительны.
Система $Pp=-g$ решается покоординатно:
\[
p_i=-\frac{g_i}{A_{ii}}.
\]
\begin{block}{Смысл}
Координата с большой диагональной кривизной получает меньший шаг. Это дёшево, но не учитывает связи между координатами.
\end{block}
\end{frame}

% 40
\begin{frame}{Пример: диагональная матрица заметно помогает}
\small
Возьмём
\[
A=\begin{pmatrix}100&9\\9&1\end{pmatrix},
\qquad P=\begin{pmatrix}100&0\\0&1\end{pmatrix}.
\]
Мы используем критерий Сильвестра для симметричной матрицы $2\times2$:
$A\succ0$ тогда и только тогда, когда $A_{11}>0$ и $\det A>0$.
Здесь $A_{11}=100>0$, $\det A=100\cdot1-9^2=19>0$.
Так мы проверили условие, нужное для оценок сходимости.
После масштабирования:
\[
B=P^{-1/2}AP^{-1/2}
=\begin{pmatrix}1&0.9\\0.9&1\end{pmatrix}.
\]
Здесь $\lambda(B)=\{0.1,1.9\}$, поэтому $\kappa_P=19$ вместо $\kappa(A)\approx535$.
\end{frame}

% 41
\begin{frame}{Какие собственные значения изменились}
\small
В исходной матрице крайние собственные значения различаются более чем в 500 раз. После диагонального масштабирования --- в 19 раз.
\figwide{eigen_before_after.pdf}
Уменьшение $\kappa_P$ объясняет ускорение на графике сходимости.
\end{frame}

% 42
\begin{frame}{Когда диагональная матрица не помогает}
\small
Пусть $A=R\operatorname{diag}(1,100)R^\top$, где $R$ --- поворот на $45^\circ$.
Тогда $\operatorname{diag}(A)=50.5I$: диагональное $P$ одинаково масштабирует обе координаты.
\figshort{rotated_failure.pdf}
Отношение собственных значений остаётся равным $100$. Причина вытянутости лежит в \textbf{повёрнутом} направлении, а не вдоль осей координат.
\end{frame}

% 43
\begin{frame}{Теорема: спектральная эквивалентность}
\footnotesize
\begin{block}{Утверждение}
Если для $A,P\succ0$ существуют $0<c\le C$ такие, что
\[
cP\preceq A\preceq CP,
\]
то все собственные значения $B=P^{-1/2}AP^{-1/2}$ лежат в $[c,C]$ и $\kappa_P\le C/c$.
\end{block}
Неравенство между матрицами означает, что для \emph{каждого} вектора $v$
\[
c\,v^\top Pv\le v^\top Av\le C\,v^\top Pv.
\]
Иными словами, $P$ оценивает квадратичную форму $A$ сверху и снизу с постоянными $c,C$.
\end{frame}

% 44
\begin{frame}{Доказательство спектральной эквивалентности}
\footnotesize
Возьмём произвольный $u\ne0$ и положим $v=P^{-1/2}u$.
Тогда $u=P^{1/2}v$, поэтому
$v^\top Pv=u^\top u$ и $v^\top Av=u^\top Bu$.
Подставляем в условие теоремы:
\[
c\,u^\top u\le u^\top Bu\le C\,u^\top u.
\]
Теперь выберем $u$ собственным вектором $B$: $Bu=\lambda u$.
После деления на $u^\top u>0$ получаем $c\le\lambda\le C$.
\par\vspace{.5em}
Так верно для каждого собственного значения, значит
$\lambda_{\min}(B)\ge c$, $\lambda_{\max}(B)\le C$ и
\[
\boxed{\kappa_P=
\frac{\lambda_{\max}(B)}{\lambda_{\min}(B)}
\le\frac{C}{c}.}
\]
\end{frame}

% 45
\begin{frame}{Вычислительная цена предобусловливания}
\small
\begin{tabular}{@{}p{.25\linewidth}p{.31\linewidth}p{.33\linewidth}@{}}
\toprule
Выбор $P$ & Применение $P^{-1}g$ & Что учитывает\\
\midrule
$I$ & без дополнительной системы & единый масштаб\\
диагональная & $O(d)$ операций & масштаб координат\\
блочно-диагональная & малые независимые системы & связи внутри блоков\\
полная $A$ & решение системы с $A$ & всю кривизну квадратичной задачи\\
\bottomrule
\end{tabular}
\par\vspace{.5em}
Хорошее $P$ уменьшает \textbf{общее} время: выигрыш в числе итераций должен окупить построение и применение матрицы.
\end{frame}

% Блок 6. Задачи ML, 9 минут. Слайды 46--51.
% 46
\begin{frame}{Предобусловливание в регуляризованном МНК}
\footnotesize
Пусть строки $X\in\R^{n\times d}$ --- объекты, столбцы --- признаки, $t\in\R^n$ --- ответы.
Ищем веса $w$, минимизируя
\[
f(w)=\frac{1}{2n}\norm{Xw-t}_2^2+\frac{\lambda}{2}\norm w_2^2,\qquad\lambda>0.
\]
Градиент и постоянный гессиан этой задачи равны
\[
\grad f(w)=\frac1nX^\top(Xw-t)+\lambda w,\qquad
A=\grad^2f(w)=\frac1nX^\top X+\lambda I\succ0.
\]
Можно взять $P=\operatorname{diag}(A)$: тогда
$P_{jj}=\norm{X_{:j}}_2^2/n+\lambda>0$.
На каждом шаге делим $j$-ю компоненту градиента на $P_{jj}$.
Так шаг по признаку с большим масштабом становится меньше, а решение системы с $P$ требует только $d$ делений.
\end{frame}

% 47
\begin{frame}{Связь с масштабированием признаков}
\footnotesize
Пусть $D=\operatorname{diag}(d_1,\ldots,d_d)$, где $d_j>0$ задаёт масштаб $j$-го признака.
Введём новые веса $u=Dw$, то есть $w=D^{-1}u$, и новые признаки $X'=XD^{-1}$.
Предсказания останутся прежними:
\[
Xw=X'u.
\]
Но штраф исходной задачи в новых весах имеет вид
$\frac{\lambda}{2}\norm{D^{-1}u}_2^2$, а не
$\frac{\lambda}{2}\norm u_2^2$.
\begin{alertblock}{Почему это важно}
Если изменить признаки и оставить прежнюю формулу штрафа для $u$, веса будут штрафоваться иначе: это уже другая оптимизационная задача.
Предобусловленный GD применяют к фиксированной $f$; он меняет шаги, а не саму функцию потерь.
\end{alertblock}
\end{frame}

% 48
\begin{frame}{Пример: разные масштабы столбцов $X$}
\small
В синтетическом МНК первый столбец признаков увеличен в 15 раз. Для гессиана с регуляризацией получаем $\kappa(A)\approx206$, а после диагонального масштабирования $\kappa_P\approx2$.
\figwide{ridge_feature_scaling.pdf}
\end{frame}

% 49
\begin{frame}{Что меняется в логистической регрессии}
\footnotesize
Для меток $t_i\in\{-1,1\}$ и строк признаков $x_i^\top$ регуляризованная потеря равна
\[
f(w)=\frac1n\sum_{i=1}^n\log(1+e^{-t_ix_i^\top w})
+\frac{\lambda}{2}\norm w_2^2,\qquad \lambda>0.
\]
Её гессиан имеет вид
\[
H(w)=\frac1nX^\top W(w)X+\lambda I,\qquad
W_{ii}(w)=\sigma_i(1-\sigma_i),\quad
\sigma_i=\frac{1}{1+e^{-t_ix_i^\top w}}.
\]
Так как $0<W_{ii}\le1/4$ и $\lambda>0$, имеем $H(w)\succ0$.
Но $\sigma_i$, а значит и кривизна, меняется вместе с $w$.
\begin{block}{Следствие для метода}
Недорогой фиксированный вариант ---
$P_{jj}=\norm{X_{:j}}_2^2/(4n)+\lambda$.
Он учитывает масштаб признаков, но не равен текущему $H(w)$; длину шага можно выбирать поиском Армихо из лекции 8.
\end{block}
\end{frame}

% 50
\begin{frame}{Как выбирать $P$ на практике}
\small
Для полезного предобусловливателя нужны три проверяемых свойства:
\begin{enumerate}
\item $P$ симметрична и положительно определена: направление убывает;
\item система $Pp=-g$ решается существенно дешевле точной системы с гессианом;
\item спектр $P^{-1/2}HP^{-1/2}$ более собран в интересующей области.
\end{enumerate}
\begin{block}{Первый кандидат}
Для задач с сильно различающимися масштабами признаков попробуйте диагональную матрицу. Затем измерьте число шагов и стоимость каждого шага.
\end{block}
\end{frame}

% 51
\begin{frame}{Ограничения и проверка результата}
\small
Диагональное $P$ не исправляет произвольный поворот собственных направлений. Плохо выбранное $P$ может даже увеличить $\kappa_P$.
Предобусловливание также не отменяет выбор длины шага и критерий остановки.
\begin{alertblock}{Что сравнивать в эксперименте}
Для одинакового начального $x_0$ сравнивайте значение $f(x_k)-f^*$ или норму градиента \textbf{и} затраченное время. Одна лишь стоимость итерации не показывает, ускорился ли алгоритм.
\end{alertblock}
\end{frame}

% Блок 7. Закрепление, 6 минут. Слайды 52--56.
% 52
\begin{frame}{Мини-задача на один шаг}
\small
Пусть $f(x)=\tfrac12(9x_1^2+x_2^2)$, $x_0=(1,2)^\top$ и $P=\operatorname{diag}(9,1)$.
\begin{enumerate}
\item Найдите $\grad f(x_0)$ и предобусловленное направление.
\item Куда приведёт шаг $\alpha=1$?
\item Какова матрица $B=P^{-1/2}AP^{-1/2}$?
\end{enumerate}
\begin{block}{Подсказка}
Сначала запишите $A$ и градиент в произвольной точке $x$.
\end{block}
\end{frame}

% 53
\begin{frame}{Мини-задача на один шаг: решение}
\small
Пусть $f(x)=\tfrac12(9x_1^2+x_2^2)$, $x_0=(1,2)^\top$ и $P=\operatorname{diag}(9,1)$.
\begin{enumerate}
\item $\nabla f(x)=(9x_1,x_2)^\top$, поэтому
$g_0=\nabla f(x_0)=(9,2)^\top$.
\item Решаем $Pp=-g_0$: $p=(-1,-2)^\top$.
При $\alpha=1$ получаем $x_{\mathrm{new}}=x_0+p=(0,0)^\top$, то есть минимум.
\item $A=\operatorname{diag}(9,1)$ и
$P^{-1/2}=\operatorname{diag}(1/3,1)$, поэтому
\[
B=P^{-1/2}AP^{-1/2}=I.
\]
\end{enumerate}
\end{frame}

% 54
\begin{frame}{Проверка понимания}
\small
\begin{enumerate}
\item Почему $P^{-1}\grad f(x)$ определяет направление убывания только при подходящем условии на $P$?
\item Какую матрицу нужно анализировать вместо $A$, чтобы оценить скорость предобусловленного GD?
\item Может ли диагональное масштабирование устранить вытянутость, направленную под $45^\circ$ к осям?
\item Почему хороший спектр ещё не гарантирует меньшее время работы?
\end{enumerate}
\end{frame}

% 55
\begin{frame}{Итоги}
\small
\[
\boxed{x_{k+1}=x_k-\alpha_kP^{-1}\grad f(x_k),\qquad
B=P^{-1/2}AP^{-1/2}.}
\]
\begin{itemize}
\item Замена координат объясняет направление предобусловленного шага.
\item При $P\succ0$ это направление убывает вне стационарной точки.
\item На квадратичной функции допустимый шаг и скорость задаёт спектр $B$.
\item Простое $P$ полезно, когда сокращение числа шагов окупает его применение.
\end{itemize}
\end{frame}

% 56
\begin{frame}{Следующая лекция: метод тяжёлого шарика}
\small
Мы научились менять геометрию градиентного шага с помощью $P$. В лекции 11 рассмотрим другой механизм ускорения: \textbf{инерцию}, которая учитывает предыдущее перемещение.
\[
\boxed{x_{k+1}=x_k-\alpha\grad f(x_k)+\beta(x_k-x_{k-1}).}
\]
На тех же квадратичных примерах исследуем, как параметры $\alpha$ и $\beta$ влияют на колебания и скорость. Это подготовит сравнение с ускоренным методом Нестерова в лекции 12.
\end{frame}

\end{document}
